\section{A positive characterization for fixed batch templates}
\label{sec:random-positive-proofs}

The logarithmic separation in Section~\ref{sec:random-schedule-separation}
does not preclude a distribution-dependent characterization on a more
restricted class. We now fix the identities and relative timing of
finitely many batches, while allowing their common time scale to have
an arbitrary finite distribution. The constants below depend on the
fixed template. In particular, the result is not a characterization
uniform in a growing number of arms or batches.

Let $B_1,\ldots,B_b$ be disjoint nonempty sets of non-anchor arms. Fix
numbers $\alpha_1>\cdots>\alpha_b=1$ and a law $Q$ supported on finitely
many positive integers. Assume $\alpha_j h$ is an integer for every
$h\in\operatorname{supp}(Q)$ and every $j$. For a known integer endpoint
$T\ge\alpha_1\max\operatorname{supp}(Q)$, draw $H\sim Q$ independently
of all reward streams and policy randomization, and set
\begin{equation}
 \tau_j(H)=T-\alpha_jH+1,\qquad i\in B_j.
 \label{eq:stretch-arrivals}
\end{equation}
Only the anchor is available before $\tau_1(H)$. Each arriving arm
remains available. Put
\[
 k_j=1+\sum_{r\le j}|B_r|,\qquad K=k_b,\qquad
 d_j=\begin{cases}\alpha_j-\alpha_{j+1},&j<b,\\1,&j=b.\end{cases}
\]
Thus the interval from batch $j$ through the round before batch $j+1$
has exactly $d_jH$ rounds. Each such length is a positive integer.
The reward laws are arbitrary independent stationary distributions on
$[0,1]$, and are fixed before $H$ is drawn. The policy knows the
template, $T$, and $Q$, and observes the arrivals causally. Let
$\mathcal Q_{T,Q,\boldsymbol\alpha,\boldsymbol B}$ denote the resulting
law of complete schedules. This information class is distinct from
that used for the anytime MC-MOSS guarantee.

Define the truncated-lifetime complexity
\begin{equation}
 W(Q)=\sup_{x\ge1}\frac{\EE_Q\min\{H,x\}}{\sqrt{x}}.
 \label{eq:stretch-complexity}
\end{equation}
It is specified directly by the arrival law, with no optimization over
bandit policies or reward environments.

\begin{theorem}[A fixed-template random-arrival characterization]
\label{thm:random-positive}
Under the preceding assumptions,
\begin{equation}
 \frac{W(Q)}{32}
 \le \cR_T^{\mathrm{joint}}
       (\mathcal Q_{T,Q,\boldsymbol\alpha,\boldsymbol B})
 \le A_{\boldsymbol\alpha,\boldsymbol B}W(Q),
 \qquad
 A_{\boldsymbol\alpha,\boldsymbol B}
 =\sum_{j=1}^b(20k_j+43)(k_j-1)\sqrt{d_j}.
 \label{eq:stretch-minimax}
\end{equation}
The upper bound is attained by the following causal rule. Play the
anchor until the first batch arrives; its arrival reveals
$H=(T-\tau_1+1)/\alpha_1$. At the start of each batch interval, reset
all empirical statistics and run fixed-horizon MOSS on the $k_j$
available arms for its known length $d_jH$.
This rule uses the endpoint and template, but need not know the
mixing probabilities in $Q$.

The complexity can be computed by a finite maximization:
\begin{equation}
 W(Q)=\max_{h\in\operatorname{supp}(Q)}
            \frac{\EE_Q\min\{H,h\}}{\sqrt h}.
 \label{eq:stretch-finite}
\end{equation}
Moreover, with
$V(Q)=\max_{h\in\operatorname{supp}(Q)}\sqrt h\,\PP_Q(H\ge h)$,
\begin{equation}
 V(Q)\le W(Q)\le2V(Q).
 \label{eq:stretch-tail}
\end{equation}
\end{theorem}

The rule in the theorem is an existence construction for a known
endpoint and a fixed template. It is not MC-MOSS, and the theorem does
not assert that retaining history is unnecessary for the original
problem. Here the regret of every interval is charged explicitly;
discarding its predecessor's observations merely gives a convenient
upper bound. The lower bound permits the policy to use all previous
observations and even the entire realized schedule.

We first establish three ingredients. The first is a deliberately
loose fixed-arm consequence of the MOSS index of
\citet{audibert2009minimax}. We prove the exact bound used here rather
than infer it from a distribution-free statement.

\begin{lemma}[A per-arm fixed-horizon bound]
\label{lem:stretch-local-gap}
Run fixed-horizon MOSS for $n\ge1$ rounds on $k\ge2$ independent
stationary $[0,1]$ arms, with all statistics initially zero. An
unobserved arm has index $+\infty$; after $r\ge1$ observations its
index is
\[
 \widehat\mu_{i,r}
       +\sqrt{\frac{\log_+(n/(kr))}{r}}.
\]
Use any fixed deterministic tie rule. If arm $i$ has gap
$\Delta_i=\mu^*-\mu_i>0$, its expected regret contribution satisfies
\begin{equation}
 \Delta_i\EE N_i(n)
 \le \min\left\{n\Delta_i,
   \frac{20k+10+16\log_+(n\Delta_i^2)}{\Delta_i}\right\}.
 \label{eq:stretch-local-gap}
\end{equation}
An optimal arm has zero regret contribution. The statement remains
valid when $n<k$, when some arms are never initialized, and when the
optimal arm is not unique.
\end{lemma}

\begin{proof}
Fix one optimal arm and put $v=n/k$. For its potential sample stream,
let
\[
 Z=\sup_{r\ge1}\left[
       \mu^*-\widehat\mu_{*,r}
       -\sqrt{\log_+(v/r)/r}\right]_+.
\]
Lemma~\ref{lem:Z} yields
$\PP(Z\ge x)\le\min\{1,5/(vx^2)\}$ for $x>0$.
Write $\Delta=\Delta_i$. On $\{Z\ge\Delta/2\}$, bounding the
contribution of arm $i$ by $n\Delta$ gives at most $20k/\Delta$ in
expectation. Its initial pull costs at most $\Delta\le1/\Delta$.
On the complementary event, choosing an observed suboptimal arm
implies that the fixed optimal arm has been observed, since otherwise
an unobserved arm would have infinite index. Index comparison gives
\[
 (\widehat\mu_{i,r}-\mu_i)_+
       +\sqrt{\log_+(v/r)/r}\ge\Delta/2.
\]
Either the positive error or the bonus is at least $\Delta/4$.
The pre-pull count $r$ is distinct at every selection of arm $i$.
Hoeffding's inequality therefore bounds the expected number of
error-type selections by
\[
 \sum_{r\ge1}e^{-r\Delta^2/8}
    =\frac1{e^{\Delta^2/8}-1}\le\frac8{\Delta^2}.
\]
This argument uses all deterministic potential prefixes and does not
condition a concentration bound on an adaptively selected count.

Suppose $n\Delta^2\ge1$. Counts below $\Delta^{-2}$ contribute at most
$\Delta^{-2}$ possibilities. If $r\ge\Delta^{-2}$ and the bonus is
at least $\Delta/4$, then
\[
 r\le\frac{16}{\Delta^2}\log\frac{n}{kr}
    \le\frac{16}{\Delta^2}\log(n\Delta^2).
\]
Thus the number of bonus-type counts is at most
$[1+16\log(n\Delta^2)]/\Delta^2$. Adding the underestimation,
initialization, error, and bonus terms gives the second expression
in~\eqref{eq:stretch-local-gap}. If $n\Delta^2<1$, the deterministic
bound $n\Delta<1/\Delta$ gives the same assertion. Combining it with
the deterministic bound $n\Delta$ proves the stated minimum.
\end{proof}

\begin{lemma}[A uniform averaging inequality]
\label{lem:stretch-average}
Let $L$ be any finitely supported positive integer random variable,
and put
$W_L=\sup_{x\ge1}\EE\min\{L,x\}/\sqrt x$.
For every $0<\Delta\le1$ and every $a,b\ge0$,
\begin{equation}
 \EE\min\left\{L\Delta,
           \frac{a+b\log_+(L\Delta^2)}{\Delta}\right\}
 \le (1+a+2b)W_L.
 \label{eq:stretch-average}
\end{equation}
If $L=dH$ for $d>0$ and both $L$ and $H$ are positive integer valued,
then $W_L=\sqrt d\,W(Q)$.
\end{lemma}

\begin{proof}
Set $x=\Delta^{-2}\ge1$. The contribution from $L\le x$ is at most
$\EE[L\ind\{L\le x\}]/\sqrt x\le W_L$. For $t\ge1$,
\[
 \PP(L>t)\le\frac{\EE\min\{L,t\}}{t}
               \le\frac{W_L}{\sqrt t}.
\]
Consequently, the contribution from $L>x$ is at most
\begin{align*}
 \sqrt x\,\EE[(a+b\log(L/x))\ind\{L>x\}]
 &=\sqrt x\left\{a\PP(L>x)
             +b\int_x^\infty\frac{\PP(L>t)}{t}\,dt\right\}\\
 &\le(a+2b)W_L.
\end{align*}
The integral identity follows from the nonnegative layer-cake
formula; finite support ensures finiteness. This proves the first
claim. The scaling identity follows from~\eqref{eq:stretch-finite},
proved below: evaluating at the support points $d h$ gives exactly
$\sqrt d$ times the expressions at $h$.
\end{proof}

\begin{lemma}[A common-environment testing bound]
\label{lem:stretch-testing}
For every $x\in\operatorname{supp}(Q)$ and every causal policy in the
stated information class,
\begin{equation}
 \sup_P\EE_{H\sim Q,P}R_T
      \ge\frac{\EE_Q\min\{H,x\}}{32\sqrt x}.
 \label{eq:stretch-testing}
\end{equation}
The inequality also holds if the policy is told the realized value
of $H$ before the experiment starts.
\end{lemma}

\begin{proof}
Choose one label $i_0\in B_b$. In two environments $P_-$ and $P_+$,
all other arms are Bernoulli with mean $1/2$, while arm $i_0$ is
Bernoulli with mean $1/2-\delta$ and $1/2+\delta$, respectively, where
$\delta=1/(8\sqrt x)\le1/8$. Both environments are fixed independently
of $H$. The equally weighted prior on them is also independent of $H$.

Condition on $H=h$. Before $i_0$ arrives, every observed reward has
the same law under the two environments. This includes all rewards
collected in the earlier batch intervals and every action based on
them. After its arrival there are exactly $h$ rounds remaining. Put
$\ell_h=\min\{h,x\}$, a positive integer, and retain the full history
through the first $\ell_h$ of those rounds. Only observations from
$i_0$ contribute to the likelihood ratio. If $N$ counts its pulls in
these $\ell_h$ rounds, the chain rule for relative entropy gives
\[
 \operatorname{KL}(\mathbb P_{-,h},\mathbb P_{+,h})
 =\EE_{-,h}N\,
    \kl(1/2-\delta,1/2+\delta)
 \le16\delta^2\ell_h\le\frac14.
\]
Here $\mathbb P_{\pm,h}$ include the entire earlier history and the
policy's randomization. The policy kernels agree under both
environments, and the schedule is the same after conditioning.
The Bernoulli inequality follows from
\[
 \kl(1/2-\delta,1/2+\delta)
 =2\delta\log\frac{1+2\delta}{1-2\delta}
 \le\frac{8\delta^2}{1-2\delta}\le16\delta^2.
\]
Pinsker's inequality yields
$\operatorname{TV}(\mathbb P_{-,h},\mathbb P_{+,h})\le1/\sqrt8<1/2$.
Because $0\le N\le\ell_h$,
$|\EE_{+,h}N-\EE_{-,h}N|\le\ell_h/2$. The regret in the retained
interval is $\delta N$ in $P_-$ and $\delta(\ell_h-N)$ in $P_+$.
Its prior average is therefore at least
\[
 \frac\delta2
       (\ell_h+\EE_{-,h}N-\EE_{+,h}N)
 \ge\frac{\delta\ell_h}{4}
 =\frac{\min\{h,x\}}{32\sqrt x}.
\]
All other regret is nonnegative. Average over $h$; one of the two
fixed environments has risk at least the prior average. This closes
the lower bound without selecting an environment after observing the
schedule. If $H$ is supplied at time zero, its conditional law and
all pre-arrival observations are still identical under $P_-$ and
$P_+$, so the proof is unchanged.
\end{proof}

\begin{proof}[Proof of Theorem~\ref{thm:random-positive}]
We first verify the finite maximization independently of the regret
argument. List the support as $h_1<\cdots<h_M$. Between consecutive
support points,
\[
 \frac{\EE\min\{H,x\}}{\sqrt x}
       =\frac{u}{\sqrt x}+v\sqrt x,
 \qquad u,v\ge0.
\]
Its derivative has the sign of $vx-u$, so any interior critical
point is a minimum. Below $h_1$ the expression is $\sqrt x$; above
$h_M$ it is $\EE H/\sqrt x$. Its maximum is consequently attained
at a support point, proving~\eqref{eq:stretch-finite}. This also
justifies the scaling assertion of Lemma~\ref{lem:stretch-average}.

For the tail comparison, $\EE\min\{H,h\}\ge h\PP(H\ge h)$ gives
$V(Q)\le W(Q)$. For every $t>0$,
$\PP(H>t)\le\min\{1,V(Q)/\sqrt t\}$: between support points use
the smallest support point strictly above $t$, and above the largest
point the tail vanishes. Hence, for $x\ge1$,
\[
 \EE\min\{H,x\}
 =\int_0^x\PP(H>t)\,dt
 \le V(Q)\int_0^x t^{-1/2}\,dt
 =2V(Q)\sqrt x.
\]
This proves~\eqref{eq:stretch-tail}.

Take a maximizing support point $x$ in
Lemma~\ref{lem:stretch-testing} to obtain the lower bound.
For the upper bound, use the rule stated in the theorem. Its
pre-first-batch regret is zero. Conditional on $H=h$, all batch
boundaries and interval lengths are deterministic, and the arriving
batch labels reveal the required schedule. In interval $j$ the
available arm set and its best mean are fixed, so each gap
$\Delta_{i,j}=\max_{r\le j,\,i'\in B_r\cup\{a\}}\mu_{i'}-\mu_i$
is the same for every value of $h$. (The maximum simply ranges over
the anchor and batches $1,\ldots,j$.) Resetting the statistics makes
this interval a fixed-horizon MOSS experiment with length $d_jh$
and $k_j$ arms.

To justify this last step in the presence of adaptive earlier
sampling, one may realize rewards as independent variables
$X_{i,t}\sim P_i$ for each label and calendar round and reveal only
$X_{I_t,t}$. This has exactly the same observed law as independent
potential pull streams. Conditional on the complete past at a batch
boundary and on $H$, future calendar variables remain independent
and keep their original laws. Thus the bound in
Lemma~\ref{lem:stretch-local-gap} applies conditionally on the past;
no old observations are treated as independent fresh data.

Apply Lemma~\ref{lem:stretch-average} with $L=d_jH$,
$a=20k_j+10$, and $b=16$ to each positive gap. Each arm contributes
at most $(20k_j+43)\sqrt{d_j}W(Q)$. At least one arm has zero gap,
so there are at most $k_j-1$ positive-gap terms. Sum over the batch
intervals. They partition every round with more than one available
arm, so all prefix and later regret has been included. This yields
the upper bound in~\eqref{eq:stretch-minimax}, uniformly over the
fixed reward laws.
\end{proof}

\begin{corollary}[When the expected deterministic capacity is correct]
\label{cor:stretch-capacity}
For a fixed template as above, write
$D(Q)=\EE_Q\sqrt H/W(Q)$. Then $D(Q)\ge1$ and
\begin{equation}
 \frac{D(Q)}{A_{\boldsymbol\alpha,\boldsymbol B}}
 \le
 \frac{\EE_Q\cC_T(\tau(H))}
      {\cR_T^{\mathrm{joint}}
       (\mathcal Q_{T,Q,\boldsymbol\alpha,\boldsymbol B})}
 \le32\sqrt{(K-1)\alpha_1}\,D(Q).
 \label{eq:stretch-capacity-ratio}
\end{equation}
Consequently, over any sequence of laws and horizons with this
template fixed, the expected deterministic capacity is of the same
order as the joint minimax risk if and only if $D(Q)$ stays bounded.
\end{corollary}

\begin{proof}
For every $h\ge1$, a single final-batch candidate has lifetime $h$,
so its feasible allocation gives $\cC_T(\tau(h))\ge\sqrt h$.
There are $K-1$ non-anchor arms, all with lifetime at most
$\alpha_1h$. Since $\phi(n)\le\sqrt n$, the total-allocation
constraint and Cauchy--Schwarz give
\[
 \cC_T(\tau(h))
 \le\sqrt{(K-1)\alpha_1h}.
\]
Average these bounds and combine them with
Theorem~\ref{thm:random-positive}. Finally,
$\min\{H,x\}/\sqrt x\le\sqrt H$ pointwise, whence
$W(Q)\le\EE\sqrt H$ and $D(Q)\ge1$.
\end{proof}

For a deterministic scale $H=h_0$, $W(Q)=\sqrt{h_0}$ and $D(Q)=1$.
If the support lies in $[h_0,Lh_0]$ for a fixed $L$, then
$D(Q)\le\sqrt L$. At the other extreme, for
$\PP(H=4^j)=c_m2^{-j}$, $0\le j\le m$,
\[
 \frac{\EE\min\{H,4^r\}}{2^r}
 =c_m\{3-2^{-r}-2^{r-m}\},\qquad0\le r\le m.
\]
Thus $1\le W(Q)<3$, whereas $\EE\sqrt H=c_m(m+1)$.
For a template with noninteger $\alpha_j$, choose a fixed positive
integer $q$ such that every $q\alpha_j$ is an integer and use
$H=q4^J$ instead; both $W(Q)$ and $\EE\sqrt H$ are multiplied by
$\sqrt q$. Such a $q$ exists for every template admitting a nonempty
legal integer support.
The same order-one versus logarithmic separation therefore persists
for every fixed multi-batch template covered by the theorem, not
only for a single arriving candidate.

The tail functional in~\eqref{eq:stretch-tail} is the familiar weak
$L^1$ tail quantity for $\sqrt H$; we do not claim a new general
probability functional. The result identifies its statistical role in this
specific independent-arrival experiment. Its limitations are material:
the template is fixed; the endpoint is known; the first arrival
reveals the remaining schedule; and the upper constant grows with
the number of arms and batches. Random schedules that do not share a
revealing common stretch, and rates uniform in growing templates,
remain outside the theorem.
